% Generated by scripts/paper/supp_experiment_record.py; do not edit.
\paragraph{SG-A3G GoCA-oriented A3 seed and paired rerun (Sec.~128).} Tests the open A3 orientation caveat directly. One diffusion run accumulates both the historical A3 seed and A3G, whose SA similarity uses GoCA's orientation; the two differ only in per-call CA module weights of resized modules, giving an exactly paired comparison per image. Each seed feeds the working default and the v2 V3 capped hierarchy. Pre-fixed questions: does G0+R2 - V0+R2 exceed the +/-0.002 tie band, and does V3's gain survive the corrected seed?

\emph{Cell names.} V0 = historical A3 $\rightarrow$ FINE-P2 (reference; same as SG-combined-v2 V0). V3 = historical A3 $\rightarrow$ v2 V3 (one-sided cap at 16$\rightarrow$32). G0 = A3G $\rightarrow$ FINE-P2. G3 = A3G $\rightarrow$ v2 V3. A3G = A3 with GoCA's affinity orientation $<$M,P$>$ instead of the historical $<$M\^{}T,P$>$. +R2 = fixed DenseCRF. Candidates V3, G0, G3.

\emph{Reference.} V0 / V0+R2 = historical A3 $\rightarrow$ FINE-P2.

\emph{Status.} Registered and submitted (protocol Sec. 128); results pending. A two-image smoke shows the seeds barely differ (mean $|$A3G - A3$|$ 0.0005 / 0.0010, Pearson r \textasciitilde{}1.0); this is an execution check, not a result. If G0+R2 - V0+R2 exceeds the tie band, A3's descendants are owed a rerun on A3G; otherwise the historical results stand with the orientation documented.

\paragraph{SG-combined-v3 single changes to the V3 capped hierarchy (Sec.~129).} Three owner-approved hypotheses, each changing exactly one thing relative to v2 V3: a per-head cap, a native coarse CA seed for the 16x16 start, and a temporal cap across SA calls. Seed is historical A3; SA calls 10-20; cap clipped mass discarded. Five cells raw and after R2, a mechanism table against V3, fixed decision rules, GT cap-loss diagnostics and an owner-approved dense-map sidecar aimed at the frame-filling recall loss.

\emph{Cell names.} R0 = FINE-P2 (= v2 V0), the registered selection reference. R1 = v2 V3 replay: r16 = head-mean 16x16 SA x c16, z32 = cap of the 32x32 head-mean pass against r16, then head-mean 64x64 pass; local mechanism reference. W1 = per-head cap: z32 = mean over heads of cap(r16, A\^{}32\_h U r16) with the head-mean parent (removes at least as much mass as R1). W2 = native coarse seed: the R1 chain starts from c16* = minmax of the A3 CA contributions at $\leq$ 16x16 modules (unchanged A3 weights), so 64-px CA no longer enters the hierarchy. W3 = temporal cap: each call's R1 map is capped by the area-downsampled mean of earlier calls' maps (first call uncapped). +R2 = fixed DenseCRF. Caution: R0/R1 are new names for v2 V0/V3.

\emph{Reference.} R0 / R0+R2 (selection reference, = v2 V0); R1 is the mechanism reference for W1-W3.

\emph{Status.} Registered and submitted (protocol Sec. 129), chained behind SG-A3G; results pending. Pre-fixed rules: W2 becomes a final-method candidate if it beats R1 in $\geq$ 4/5 folds with the paired interval above zero and VOC sheep does not fall; W1 closes if sheep falls vs R1; W3 closes unless it beats R1 in $\geq$ 4/5 folds with the interval above zero. Caveat: A3's SA similarity uses a transposed affinity orientation relative to GoCA (protocol Sec. 107); on four images the correction changes per-call module weights by 0.9-2.0\% (Sec. 114), and a corrected replay has not been run.

